\section{Scope}
\label{sec:intro:scope}

The scope of this thesis is bounded along four axes: the data
sources analysed, the analytical questions addressed, the
stakeholders supported, and the deliverables produced.

\paragraph{Data scope.}
The analysis covers two complementary HCMC bus-network data streams:
GPS positioning data (a single-day local sample of 720{,}088 raw
pings from 2025-03-18 covering 9 routes; the full 50-day dataset of
$\sim$34\,GB referenced in the dataset specification~\cite{hcmcgps2025}
is reachable but not loaded locally), and AFC ticketing data
(17.9~million boarding events with paired payments over the 90-day
window 2026-01-01 to 2026-03-31). Static reference data --- stops,
routes, vehicles, companies --- is included as dimension tables.
The two streams are not joined by timestamp: their calendar windows
do not overlap. They are therefore treated as complementary sources
and combined narratively in the dashboard rather than fused at the
row level.

\paragraph{Analytical scope.}
Six classes of analytical question are in scope and one class is
explicitly out of scope. In scope:
(i)~system-wide demand and reliability (\emph{Network Health \&
Performance});
(ii)~when and where people ride (\emph{Demand \& Crowding});
(iii)~the network as a city map, including the bus--metro feeder
matrix (\emph{Spatial \& Connectivity});
(iv)~route profit-and-loss, payment-channel adoption, and
promotion impact (\emph{Financial \& Operations});
(v)~per-route deep-dive and cross-route clustering
(\emph{Routes}); and (vi)~the flagship lens on how universities
and schools shape demand (\emph{School-Aware Demand}). Out of
scope: real-time monitoring (the dashboard is historical-only),
individual passenger tracking (data killed it: 0.1\% of tickets
populate \texttt{user\_id}), passenger satisfaction (no survey
data exist), and natural-language SQL (delivered by an external
team).

\paragraph{Stakeholder scope.}
The dashboard is designed for four user roles, used internally as
a feature-prioritisation tool but not exposed in the UI:
\emph{passenger} (crowding, reliability, beginner-friendliness),
\emph{driver} (own shift, route detail),
\emph{company} (fleet performance, revenue, scorecard versus
network), and \emph{admin / MCPT} (network performance,
anomalies, equity, capacity). Each card is anchored internally to
one or two roles but exposes the same chrome to every user; no
role switcher, no role badges. The role-aware design is
implemented as feature prioritisation rather than separate views
so a stakeholder who needs cross-card context never has to leave
the tool.

\paragraph{Deliverables.}
Phase~1 produces four artefacts: (i)~this written report;
(ii)~four executed exploratory analysis notebooks
(\texttt{notebooks/01--04}) that ground every design decision in
the report; (iii)~the high-fidelity UI mockup at
\texttt{docs/design/app-layouts.html} that locks the dashboard's
visual language; and (iv)~the data-pipeline skeleton at
\texttt{src/duckdb\_loader.py} that already produces the
twenty-four raw views the warehouse design extends. The full
Streamlit application and the eight pre-computed aggregate views
are Phase~2 deliverables, scheduled per the implementation
roadmap in Chapter~\ref{chapter:future}.
